\documentclass[12pt]{article}
\usepackage{graphicx} % Required for inserting images
\usepackage[utf8]{inputenc} % Codificación de caracteres
\usepackage[spanish]{babel} % Idioma del documento
\usepackage{amsmath, amssymb, amsthm}
\usepackage{pgfplots}
\usepackage{tabularx}
\usepackage[margin=1in]{geometry}
\usepackage{svg}
\usepackage{float}
\usepackage{fancyhdr}
\usepackage{setspace}
\usepackage[most]{tcolorbox}

\newtcolorbox{tadelis}{
    colback=blue!5,
    colframe=blue!60!black,
    title=\textit{Game Theory, Steven Tadelis},
    fonttitle=\bfseries,
    boxrule=0.8pt,
    arc=2mm
}
\newtcolorbox{comentarioprofe}{
    colback=green!5,
    colframe=green!50!black,
    title=Comentario,
    fonttitle=\bfseries,
    boxrule=0.8pt,
    arc=2mm
}
\newtcolorbox{importante}{
    colback=orange!8,
    colframe=orange!70!black,
    title=Importante,
    fonttitle=\bfseries,
    boxrule=0.8pt,
    arc=2mm
}
\newtcolorbox{ejemplo}[1]{
    colback=gray!8,
    colframe=gray!70!black,
    title=#1,
    fonttitle=\bfseries,
    boxrule=0.8pt,
    arc=2mm
}

\geometry{margin=1in}
\onehalfspacing
\pagestyle{fancy}
\rhead{Teoría de Juegos e Información}
\lhead{Incertidumbre y Decisiones Racionales}
\title{Tema 1. Incertidumbre y Decisiones Racionales}
\author{Ignacio Sanabria}
\date{Agosto 26}
\begin{document}
\maketitle

\section{¿Que es la teoría de juegos?}
La teoría de juegos busca explicar la forma en que se toman decisiones racionales en el día a día. \\
\textbf{Agente Racional:} Un agente racional es aquel que busca siempre obtener su mayor utilidad. \\
\subsection{Conceptos Básicos}
\begin{itemize}
    \item Triviales: no importan tanto, son superficiales.
    \item Trascendentales: Afectan mis pagos esperados a futuro. 
    \item \textbf{Supuestos:} vitales para explicar un modelo. Un exceso de supuestos convierte el modelo en inexacto. 
\end{itemize}

\section{Básicos de Incertidumbre}
Se platea un problema de decisión que incluye: 
\begin{itemize}
    \item Situaciones sobre las que se debe decidir. 
    \item Conjunto de agentes económicos: Jugadores $J=\{\}$
    \item Conjunto de acciones a tomar: Acciones $A=\{\}$
    \item Conjunto de resultados: Resultados $R=\{\}$
\end{itemize}
Para que un problema de decisión tenga sentido, los agentes económicos deben tener preferencias sobre los resultados. \\
Estas vienen dadas por: 
\begin{enumerate}
    \item $x\succeq y \xrightarrow{}$ Débilmente Preferido 
    \item $x\succ y \xrightarrow{}$ Estrictamente preferido
    \item $x\sim y \xrightarrow{}$ Indiferente
\end{enumerate}

\section{Axiomas de Preferencias}
\begin{enumerate}
    \item \textbf{Axioma de Completitud:} Los agentes económicos siempre son capaces de ordenar sobre dos resultados. \textit{No pueden quedar indiferentes ante dos resultados.}
    \item \textbf{Axioma de Transitividad:} Son capaces de ordenar todos los resultados. \textit{Garantiza que no existen contradicción de elecciones de los jugadores.} 
\end{enumerate}
\subsection{Paradoja de Condorcet 1785}
Condorcet planteo que es posible al tener un grupo de agentes individualmente racionales, y ponerlos a tomar una decisión en grupo, estos empiecen a tomar decisiones irracionales como grupo. \\
Sea $J=\{A,B,C\}$ un conjunto de estudiantes que deben elegir si prefieren un quiz teórico ($t$), practico ($p$), o mixto ($m$). 
\begin{align*}
    \text{Con preferencias:}&\\\\
    A:& \quad t\succ m\succ p \\
    B: & \quad m \succ p \succ t \\
    C: & \quad p \succ t \succ m 
\end{align*}
Si les pedimos votar obtendriamos: 
\begin{align*}
    \text{t vs m:}& \quad t \succ m \quad \textit{A favor de t: A,C}\\
    \text{p vs m:}& \quad m \succ p \quad \textit{A favor de m: A,B}\\
    \text{t vs p:}& \quad p \succ t \quad \textit{A favor de p: B,C}\\
\end{align*}
Sin embargo, esto rompe el \textbf{axioma de transitividad} dado que:
\begin{align*}
   \text{Por transitividad:}& \quad  t \succ m \succ p  \\
   \text{pero la ultima votacion coloca:}&\quad  p \succ t \\
   \therefore t \succ p \;\mathbf{\Rightarrow \Leftarrow} \;p \succ t
\end{align*}

\section{Decisiones}
A lo largo del curso las decisiones se verán representadas en arboles de decisión, tal cual se muestra continuación: 
\begin{figure} [H]
    \centering
    \includesvg[width=0.5\linewidth]{figuras/juegos_1.drawio.svg}
    \caption{Árbol de Decisión}
    \label{fig:placeholder}
\end{figure}
\subsection{Paradigma de la elección racional}
Elegir la mejor acción del conjuntos $A$ dadas su preferencias y restricciones. 
Tal que parten de $A\rightarrow X \Rightarrow$ Preferencias al respecto. \\
Para esto se establece una función de pago que viene a representar la asociación de preferencias de la elección. \\
La función de pagos esta dada por: $$v(a)=u(x(a))$$

\section{Incertidumbre}
Los resultados estocásticos hacen referencia a la incertidumbre asociada a cada resultado. \\
\textbf{Supuestos}
\begin{enumerate}
    \item Cada resultado se puede determinar objetivamente. 
    \item Los tomadores de decisiones se interesan en el resultado, no en el proceso para llegar a el. 
    \item $X\in R$, siendo $R$ el conjunto finito de resultados. 
    \item cada posible resultado esta presente en la probabilidad. 
\end{enumerate}
A continuación iremos trabajando un ejemplo con diversos conceptos relacionados. 
\begin{ejemplo}{Ejemplo Proyecto de Investigación y Desarrollo Parte 1}
    Una empresa $E$ debe elegir si realizar ($g$) o no realizar ($s$) un proyecto de investigación para mejorar su negocio. Ademas saben que en sus resultados $X$ obtienen un $10$ si resulta exitoso, y $0$ si no. \\
    Por lo tanto su conjunto de acciones ($A$) y resultados ($X$) es:
    $$A=\{g,s\}\\X=\{0,10\}$$\\
    Ademas se sabe que si deciden realizarlo $g$ tienen un 10\% de probabilidad de obtener $10$, y si deciden no realizarlo $s$ tienen un 50\% de obtener 10. \\
    Un árbol de decisiones se vería de la siguiente manera: \\
    \begin{figure} [H]
        \centering
        \includesvg[width=0.5\linewidth]{figuras/juegos_2.drawio.svg}
        \label{fig:placeholder}
    \end{figure}
\end{ejemplo}
\subsection{Lotería Simple}
Una lotería simple esta dada por un conjunto de resultados $X=\{x_1,x_2,...,x_n\}$, se define una probabilidad  $P=\{p(x_1),p(x_2),...,p(x_n)\}$ tal que $$\sum_{k=1}^np(x_k)=1$$
La importancia de las loterías es marcar la presencia de la naturaleza en la toma de decisiones, estas al trabajar con probabilidades se basan en ideas básicas de las distribuciones de probabilidad. 
\subsubsection{Teorema de Bayes}
Sea $\{A_1,A_2,...,A_n\}$ un conjunto de eventos mutuamente excluyentes y exhaustivos, donde $P(A_i)>0$. Y sea B un suceso del que se conocen sus probabilidades condicionales $P(B|A_i)$. \\
Entonces el teorema de Bayes muestra que: 
$$P(A_i|B)=\frac{P(B|A_i)P(A)}{P(B)}$$
\subsection{Lotería Compuesta}
Cuando se tiene una lotería sobre una lotería, se clasifica como lotería compuesta. 
\begin{ejemplo}{Ejemplo Proyecto de Investigación y Desarrollo. Parte 2}
    En este caso estamos cargando de probabilidad a la segunda lotería, podemos ver que la parte donde la naturaleza ($N$), ahora parte de nuevas probabilidades. 
    \begin{figure} [H]
        \centering
        \includesvg[width=0.4\linewidth]{figuras/juegos_3.drawio.svg}
        \label{fig:placeholder}
    \end{figure}
\end{ejemplo}
\subsubsection{Lotería Simple Continua}
Sobre el intervalo $X=\{ {{\underline{x}},\bar{x}}\}$ que viene de una Función de Distribución Acumulada (FDA) $F:X\rightarrow[0,1]$ donde $F(\hat{x})=P(x \leq \hat{x})$ es la probabilidad de que un resultado sea menor o igual a $\hat{x}$. 
\begin{ejemplo}{Ejemplo Proyecto de Investigacion y Desarrollo. Parte 3}
    Ahora asumamos que existe un costo ($c$) por realizar ($g$) el proyecto de investigacion, de forma que nuestro arbol de decisiones se veria asi. 
     \begin{figure} [H]
        \centering
        \includesvg[width=0.4\linewidth]{figuras/juegos_4.drawio.svg}
        \label{fig:placeholder}
    \end{figure}
    En este caso podemos ver la utilidad esperada para determinar cual es el $C$ que pueda tomar la empresa. 
    \begin{align*}
        \pi_s&=0.5\cdot 10+0.5\cdot 0=5\\\\
        \pi_g&=0.625[0.9\cdot(10-c)+0.1(0-c)]+0.375[0.5(10-c)+0.1(0-c)]\\
        \pi_g&=7.5-c\\
        & \text{Para ser congruentes se debe cumplir que:} \\ \pi_g&>\pi_s \\ 7.5-c&>5 \\ 2.5&>c 
    \end{align*}
    Por lo tanto descubrimos que el costo minimo para decidir hacer o no la investigacion es $2.5$. 
\end{ejemplo}

\section{Axioma de Independencia}
Se consideran cuatro loterías, tal que: 
\begin{itemize}
    \item Lotería I: resultados R con probabilidad de $1-\alpha$ y P con probabilidad Q. 
    \item Lotería II: resultados R con probabilidad de $1-\alpha$ y P con probabilidad $\alpha$.
    \item Lotería III: resultados S con probabilidad de $1-\beta$ y P con probabilidad $\beta$. 
    \item Lotería IV: resultados S con probabilidad de $1-\beta$ y Q con probabilidad $\beta$. 
\end{itemize}
Establece que: 
\begin{align*}
    L_I>L_{II} \iff P>Q \\ 
    L_{III}>L_{IV} \iff Q>P
\end{align*}
\newpage
\section{Riesgo}
Para la seccion de riesgo vamos a iniciar con un ejemplo: 
\begin{ejemplo}{Ejemplo de Riesgos}
    Sea $X_1=\$4,\;X_2=\$9,\; X_3=\$16,$ un conjunto de acciobes para el jugador $J=\{Y\}$. \\
    Para estas acciones se jugaran dos loterias con las siguientes probabilidades: \\
    $$P_1=(\frac{7}{12},0,\frac{5}{12}), \quad P_2=(0,1,0)$$
    \begin{enumerate}
        \item {Ordene las preferencias monetarias}
        $16\succ9\succ4$
        \item {Ordenamiento de la Utilidad}
        $u(16)\succ u(9)\succ u(4)$
        \item Considere las siguientes funciones de utilidad y calcule el pago esperado de cada loteria con ellas. 
        \begin{itemize}
            \item $X=x$
            \begin{align*}
                \frac{7}{12}\cdot x_1 + \frac{5}{12} \cdot x_3 = x_2 \\ 
                \frac{7}{12}\cdot 4 + \frac{5}{12} \cdot 16 = 9 \\
            9=9 \; \Rightarrow\textbf{Neutral al riesgo}
            \end{align*}
            \item $X=\sqrt{x}$
            \begin{align*}
                \frac{7}{12}\cdot \sqrt{x_1} + \frac{5}{12} \cdot \sqrt{x_3} = \sqrt{x_2} \\
                \frac{7}{12}\cdot \sqrt{4} + \frac{5}{12} \cdot \sqrt{16} = \sqrt{9} \\
                \frac{17}{6}\sim2.8333\bar{3} \prec 3 \; \Rightarrow \textbf{Aversa al riesgo}
            \end{align*}
            \item $X=x^2$
            \begin{align*}
                \frac{7}{12}\cdot x_1^2 + \frac{5}{12} \cdot x_3^2 = x_2^2 \\ 
                \frac{7}{12}\cdot 4^2 + \frac{5}{12} \cdot 16^2 = 9^2 \\
                116\succ81\; \Rightarrow \textbf{Amante al riesgo}
            \end{align*}
        \end{itemize}
    \end{enumerate}
\end{ejemplo}

\subsection{Supuestos de decisiones racionales bajo incertidumbre}
\begin{itemize}
    \item Las elecciones de acciones implican elegir loterias. 
    \item El jugador conocer la probabilidad asociada a cada loteria. 
    \item El jugador sabe que los resultadls estan condicionados a las acciones. 
\end{itemize}

\begin{ejemplo}{Estudiar o no un posgrado}
Tome los siguientes datos: 
    \begin{center}
            \begin{tabular}{|c|c|c|c|c|}
                \hline
                "" & Posgrado & Prob. Posgrado & No Posgrado & Prob. No Posgrado \\
                \hline
                Ingreso Max & 32 & 0.25 & 12 & 0.25  \\ \hline
                Ingreso Med & 16 & 0.50 & 8 & 0.50 \\ \hline
                Ingreso Min & 12 & 0.25 & 4 & 0.25 \\ \hline
                Costo       & 10 &  & 0 &  \\ 
                \hline
            \end{tabular}
    \end{center}
En base a esta informacion considere cual es la mejor opciones y cuales son sus pagos esperados: \\
Sea $X_1$ Posgrado y sea $X_2$ No Posgrado, ya con la reduccion del costo, tal que: \\ 
$X_1=\{2,6,22\},\; X_2=\{4,8,12\}$\\ 
Y sean las probabilidades dadas por: $P_1=(0.25,0.5,0,25)=P_2$\\
El arbol de decicion se veria de la siguiente forma: 
    \begin{figure} [H]
        \centering
        \includesvg[width=0.3\linewidth]{figuras/juegos_5.drawio.svg}
        \label{fig:placeholder}
    \end{figure}
Calculando el valor esperado de cada uno obtenemos: 
\begin{align*}
    E(P)=0.25 \cdot 22 + 0.5 \cdot 6 + 0.25 \cdot 2= 9 \\
    E(Np)=0.25 \cdot 12 + 0.5 \cdot 8 + 0.25 \cdot 4= 8 \\
\end{align*}
Por ende, la decision racional es decidir estudiar el posgrado. 
\end{ejemplo}
\begin{importante}
    Recuerde que una utilidad esperada no es lo mismo que un pago esperado. La utilidad tiene valor cardinal y monetario segun la unidad de medida, mientras que el pago esperado es la "esperanza" de las probabilidades. 
\end{importante}
\subsection{Dinero en el Tiempo}
\begin{itemize}
    \item \textbf{Valor del dinero en el tiempo} \\ 
    El valor presente del dinero se expresa: $\frac{x}{(1+r)^t}$
    \item \textbf{Factor de descuento}\\
    $\delta\in(0,1)$ posibilidad de que los pagos concretarse. 
    \item \textbf{Valor descontado de pago a futuro}\\
    $$\sum_{t=1}^T \delta^{t-1}u(x_t)$$ 
\end{itemize}


\begin{tadelis}
Según el libro, la ecuación de Euler caracteriza
la senda óptima de consumo.
\end{tadelis}
\begin{comentarioprofe}
El profesor enfatizó que esta condición solamente se cumple
bajo determinadas condiciones de interioridad.
\end{comentarioprofe}
\begin{importante}
La tasa de interés afecta el costo de oportunidad
del consumo presente..
\end{importante}

\end{document}